\section{CSS memory circuits and detector evidence}

\subsection{Initialization, extraction, and boundaries}
A finite CSS presentation does not determine its physical measurement circuit.
The circuit layer adds ordered gates, explicit faults, measurement-record
indices, and detector/observable parities. The reference memory schedule
initializes a known logical X or Z eigenstate: independent code stabilizers
and the chosen logical generators form a complete positive stabilizer set.
This is an ideal state boundary, not an arbitrary-input encoder or a
fault-tolerant preparation procedure.

Each retained check has a dedicated ancilla reset before every use. A Z
check uses data-to-ancilla CX gates followed by ancilla Z measurement.
An X check surrounds ancilla-to-data CX gates with ancilla Hadamards before
Z measurement. Checks run sequentially, retaining redundant and zero-weight
rows. An optional undirected interaction graph constrains these extraction
edges; it does not constrain the ideal preparation synthesis.

Writing $m_{r,j}$ for check $j$ in round $r$, the initial and temporal
parities are respectively $m_{0,j}$ and $m_{r,j}+m_{r-1,j}$ over $\F$.
Final destructive readout closes only checks matching the memory basis:
\begin{equation}
 m_{T-1,j}+\sum_{q\in\operatorname{supp}(H_{b,j})}m_q^{\rm final}.
\end{equation}
The same-basis logical representatives define final observable parities.
Opposite logical bases require separate experiments. For $m$ retained
checks and $m_b$ matching-basis checks, $T$ rounds produce $Tm+n$
measurements and $Tm+m_b$ detectors. An empty logical-observable set is
not accepted as evidence of memory protection.

\subsection{Fault and detector semantics}
The current fault model applies data depolarization at each round's start,
reset flips after ancilla resets, one-qubit depolarization after Hadamards,
and two-qubit depolarization after CX gates. X before Z measurement and Z
before X measurement model readout flips for the discarded/reset measured
qubits in this schedule. Ideal initialization has no faults. Round ticks
supply neither calibrated idle durations nor optimized gate depth.

The detector model stores sparse binary maps $H$ and $L$ on independent
Bernoulli fault mechanisms $f$:
\begin{equation}
 d=Hf,\qquad \ell=Lf.
\end{equation}
A column can touch arbitrarily many detectors or only logical observables;
there is no graphlike assumption. A decoder predicts $\hat f$ satisfying
$H\hat f=d$, but logical success additionally requires $L\hat f=\ell$.
Independent syndrome verification precedes logical comparison. Native
model samples and direct circuit samples carry distinct provenance.

\subsection{Optional compilation and limits}
The Stim backend compiles the ideal signed preparation and checks its
synthesized tableau exactly before appending extraction and record
annotations~\cite{gidney2021,stim116}. Strict detector conversion forbids random
gauge detectors, disjoint-error approximations, and automatic graphlike
decomposition. Generic signed preparations use their actual noiseless
reference; the compiler does not force all reference parities to zero.

Representation sizes, preparation/reference tableaux, synthesized output,
model expansion, decoding graphs, and sample buffers have separate guards.
Foreign synthesis and simulator workspace are not total-memory guarantees.
Reproducibility records retain circuit/model descriptions, source and
backend versions, noise placement, seeds, and batching. Seeded Stim output
is not guaranteed bitwise identical across versions or machine capabilities.

These are circuit-level memory experiments under stated assumptions.
Check sparsity does not certify fault tolerance of the extraction gadgets,
static code distance does not equal circuit distance automatically, and
finite-size failure curves alone establish no threshold.
